\documentclass{article}


\title{Chem 001A\\ Week 9}
\author{Danny Topete}
\date{November 24, 2025}

\begin{document}
\maketitle{}

\section{Monday Lecture: Nov 24}
\subsection{Gas-Evolving Reactions}
\begin{itemize}
  \item If a reaction creates a gas, then it is gas evolving
  \item Gas evolving if it is directly produced or through subsequent composition
    of a reaction
  \item Only gas evolving if gas is on the right and only the right as a product
  \item Creates gas from liquid, solid, or aquous
\end{itemize}

\subsection{Acid Base Chemistry}
\begin{itemize}
  \item Covering common acids and bases
  \item \textbf{Acid base reaction:} A reaction in which a proton is moving from one species to another.

    Reaction is one in which hydrogen ion $(H^+)$ is transferred from one chemical species to another.
  \item Acidic solution causes the proteins to denature, the unfolding process.
  \item Heartburn, the solution is an anti-acid. You need to eat a strong base. It will neutralize the HCl stomach
    acid. Creating water in your stomach.

    Acid in your stomach is important to absorb nutrients and prevent pathogens.
  \item \textbf{Acid:} any substance that produces $H^+$ or $H_3O^+ $ ions in solution
  \item \textbf{Base:} substance that produces $OH^-$ ions in aqueous solution
  \item When you \textbf{neutralize reaction}, you combine a base and an acid and create $H_2O$
    It is typically a water and a salt. as an outcome
  \item Did a proton transfer? It is typically a water and a salt.
  \item acid-base reaction,
  \item Compounds that lost a proton are an acid, ones that gain a proton are a base
  \item \textbf{Important memorize box about strong acids and strong base.}

    Strong acids have H in the beginning, strong bases end in OH
  \item Strong means they give up all of their $H^+$
  \item Weak acids, hydrofloric HF, Formic $HCHO_2$ and acetic acids $HC_2H_3O_2$
\end{itemize}

\subsection{Oxidation reduction reactions}
\begin{itemize}
  \item Last class of reactions we will cover
  \item \textbf{Example:} Breaking down glucose to create energy, or macronutrients.

    Breaking down any proteins to make energy is oxidation reduction
  \item \textbf{Redox Reactions: Oxidation-reduction reactions:}
  \item Iron rusts, when you go to oxygen rich reaction, it forms an iron oxide (rust)

    Combustion of hyydrogen, combustion of octane,

    forming table salt (forming ionic compound, electrons moved), oxygen is not present,
    but the criteria doesn't mean that oxyfen is required
  \item Redox reaction gives more bonds to oxygen
  \item How do we identify them without the presence of oxygen?
  \item \textbf{Oxidation Number:} Helps categorize oxidation reactions. The charge the atom
    would take if all shared electrons in a compound were assigned to them atom with the
    greatest attraction for those electrons.

    The electronegative takes all the electrons. Another fake calculation.
  \item Oxidation states do -1 and +1 meanwhile charges are 2+ and 2-.
    Sign goes before magnitude
  \item Oxidation state of an atom in a free elemtn is 0

    $Cu $, 0 ox state; $Cl_2$, 0 ox state.
  \item Oxidation stata of monoatomic ion is equal to its charge.

    $Ca^{2+}, +2 $ ox state.

  \item Sum of oxidation states of all atoms in
    \begin{itemize}
      \item neutral molecule or formula is 0
      \item an ion equal to the charge of the ion
      \item $H_2O$, 2(H ox state) +1(O ox state) = 0
      \item $ NO_3^-$
    \end{itemize}
  \item Group 1A metals are always ox state +1

  Group 2A metals always have an ox state of +2

  $CaF_2 $, +2 ox state
  \item Alkaline and Alkaline earth metals are always automatically assigned the +1 and +2 perspectively
  \item Rule 3 must always hold.
  \item (There are 5 rules to this...)

  \item This requires lots of practice to figure out oxidation states
    and evaluate if we have a redox reactions
  \item There must be a change in oxidation reaction in order to have a redox reaction
  \item C + 2S \rightarrow{} $CS_2$

    0 + 0 \rightarrow{} +4 -2

  \item \textbf{Oxidations} increasing in oxidation number. Electrons are being lost
  \item \textbf{Reduction} decrease in oxidation number. Electrons are being gained
  \item You can simulatnously have oxidation and reduction reactions

\end{itemize}


\end{document}

